\documentclass[11pt]{article}
\usepackage[margin=0.9in]{geometry}
\usepackage{amsmath,amssymb,amsthm,booktabs}
\usepackage[hidelinks]{hyperref}
\usepackage{xurl}
\newtheorem{theorem}{Theorem}
\newtheorem{lemma}[theorem]{Lemma}
\newtheorem{proposition}[theorem]{Proposition}
\newcommand{\rk}{\operatorname{rank}}
\newcommand{\Isom}{\operatorname{Isom}}
\newcommand{\Out}{\operatorname{Out}}
\newcommand{\PSL}{\operatorname{PSL}}
\title{Ordinary knot-group and branched-cover ranks from hyperbolic symmetry}
\author{}
\date{18 November 2026}
\begin{document}\maketitle
\begin{abstract}
Thirty-two named knots of meridional rank four have ordinary group
rank exactly three. For an explicit prime four-bridge twist family,
the double branched-cover groups have rank three and both the knot
and order-two orbifold groups have rank between three and four.
The rank bounds use the inversion of a generating pair in a Kleinian
group, together with verified canonical triangulations and the
symmetry restriction from a short filling core. At the base parameter,
two geodesics pass the short-length cutoff. Their homology classes
prevent an inversion of a generating pair from exchanging them.
\end{abstract}

Let $K_n$ be the closure of
\[
\sigma_1\sigma_2\sigma_3^{-2}\sigma_2\sigma_3^{-2}\sigma_2
\sigma_1^2\sigma_2^3\sigma_3^{-1}\sigma_2^{3+6n},\qquad n\ge0.
\]
Put $X_n=S^3\setminus K_n$, $G_n=\pi_1(X_n)$, and let
$\Gamma_n=G_n/\langle\!\langle\mu^2\rangle\!\rangle$ be the group of the
order-two knot orbifold. Its index-two parity kernel $\Pi_n$ is the
fundamental group of the double branched cover $Y_n=\Sigma_2(K_n)$.
The companion proofs give $\operatorname{mr}(K_n)=b(K_n)=4$,
the actual marked Dehn filling descriptions, and the hyperbolic
structures of $Y_n$ and $(S^3,K_n,2)$ \cite{family,covers,scope}.
Every nonidentity torsion element of $\Gamma_n$ is an involution.
The least number of torsion generators of $\Gamma_n$ is four:
all these involutions are meridian conjugates and their fixed $A_8$
quotient has double-transposition class rank four \cite{scope}.

\begin{theorem}\label{thm:family}
For every $n\ge0$,
\[
        \rk(\Pi_n)=3,\qquad 3\le\rk(\Gamma_n)\le\rk(G_n)\le4.
\]
For every $n\ge1$, the complete hyperbolic exterior $X_n$ has trivial
isometry group.
\end{theorem}

\begin{theorem}\label{thm:named}
Each of the thirty-two knots in the three-generator presentation
certificate \cite{uppers} has
\[
          \rk\pi_1(S^3\setminus K)=3<4=\operatorname{mr}(K)=b(K).
\]
These are the knots listed in Section~\ref{sec:named}.
\end{theorem}

The lower bound comes from a constraint on two generators. A
hyperbolic lattice generated by two elements has a projective
involution that inverts both. If the lattice is torsion-free then
the resulting manifold isometry is nontrivial. In the knot orbifold,
an involution inside the lattice would instead give three torsion
generators. We use the branched cover to exclude the other possibility.

\section{Inversion and marked fillings}
\begin{lemma}\label{lem:inversion}
Let $\Lambda\le\PSL_2(\mathbb C)$ be a finite-covolume lattice
with a generating pair $(a,b)$. There is a projective involution
$J\ne1$ normalizing $\Lambda$ and conjugating $a,b$ to their inverses.
\end{lemma}
\begin{proof}
This is the generating-pair involution of J\o rgensen, as used in
\cite[Lemma~22]{bw}. Here is a matrix proof valid also for lattices
with torsion. Lift $a,b$ to $A,B\in\mathrm{SL}_2(\mathbb C)$.
Write $A_0=A-\tfrac12\operatorname{tr}(A)I$ and likewise $B_0$.
Choose a nonzero trace-zero matrix $C$ satisfying
$\operatorname{tr}(CA_0)=\operatorname{tr}(CB_0)=0$. Such a matrix
exists because there are two linear equations in a three-dimensional
space. If $C$ were singular, it would be conjugate to a nonzero
strictly upper triangular matrix. The two trace equations would
then force $A,B$ to be upper triangular in that basis. A finite-volume
hyperbolic lattice is non-elementary and fixes no common ideal point.
Thus $C$ is invertible. For trace-zero matrices $U,V$ the sum
$UV+VU$ equals $\operatorname{tr}(UV)I$, by the quadratic
Cayley--Hamilton identity applied to $U,V,U+V$. Therefore $C$
anticommutes with $A_0,B_0$. This gives $CAC^{-1}=A^{-1}$ and
$CBC^{-1}=B^{-1}$. The projective class of $C$ has order two since
$C^2$ is scalar and $C$ itself is not scalar.
\end{proof}

\begin{lemma}\label{lem:drilling}
Let $M$ be a complete finite-volume hyperbolic manifold and $\gamma$
an embedded geodesic. Suppose $M\setminus\gamma$ has a complete
finite-volume hyperbolic structure $P$. The subgroup of
$\Isom(M)$ preserving $\gamma$ has order at most $|\Isom(P)|$.
\end{lemma}
\begin{proof}
Restriction defines a homomorphism from the stabilizer to
$\Out(\pi_1(P))$, identified with $\Isom(P)$ by Mostow--Prasad
rigidity \cite[Theorem~1.10]{afw}. If the restricted outer action
is trivial, its action on $\pi_1(M)$ is inner, because the filling
map $\pi_1(P)\to\pi_1(M)$ is surjective. Rigidity then makes the
original isometry the identity. This proves injectivity, without
requiring the restriction of the filled metric to be complete.
\end{proof}

Let $P$ be the two-cusped exterior given by adding the braid suffix
$\sigma_4\sigma_3\sigma_2^2\sigma_3\sigma_4^{-1}$ to the literal
seventeen-crossing four-braid $K_0$. The fifth strand closes to the
crossing circle. In the preferred meridian-longitude coordinates,
the signed disk-twist argument in \cite{covers} gives $X_n$ by
$(1,-3n)$ filling of the crossing-circle cusp. In the retained
triangulation this is cusp zero. Cusp one is the knot cusp and
remains unfilled. The same marked covering argument in \cite{covers}
gives $Y_n=N(-1,2+6n)$ for a fixed one-cusped, fourteen-tetrahedron
parent $N$. We use these identifications of the actual pairs.

\section{Verified geometry and geodesics}\label{sec:geometry}
\begin{proposition}\label{prop:geocert}
The complete hyperbolic structures on $P,N$ satisfy the following
bounds. $L_P(n)$ and $L_N(n)$ are the normalized lengths of the
specified filling slopes.
\begin{center}
\begin{tabular}{lccc}
\toprule
Parent & $|\Isom|$ upper bound & Systole lower bound & Normalized length lower bound\\
\midrule
$P$ & $1$ & $1/2$ & $1+\frac52 n$\\[2pt]
$N$ & $2$ & $1/2$ & $3+5n$\\
\bottomrule
\end{tabular}
\end{center}
All geometric lower bounds in this table are strict. The three
exteriors $X_1,X_2,X_3$ have no nonidentity isometries. In the filled
metric on $Y_1$ only the marked filling core has length below $1/10$.
For $Y_0$, exactly two geodesics have length below $29/50$. One is
the marked core and the two homology classes in $\mathbb Z/117$ are
$-61$ and $43$, respectively, up to orienting the geodesics.
\end{proposition}
\begin{proof}
The companion computation uses Sage~10.9 and SnapPy~3.3.2 with
verified interval arithmetic \cite{snappy}. For $P$ it regenerates
Spherogram's 23-crossing augmented diagram. The unknotted component
is component zero, and the retained triangulation agrees exactly
in gluing and peripheral data with the freshly constructed exterior.
The $N$ triangulation is the marked parent from the cover proof.

For each cusped manifold in the certificate, a geometric canonical
triangulation is certified by positive shapes, logarithmic gluing and
completeness equations, and strictly negative face tilts. For $P$ the
cusp areas are equalized before testing the tilts. The strict tilt
criterion from the Epstein--Penner construction and Weeks's tilt
formula ensures that the canonical cells are exactly the tetrahedra
\cite[Sections~3.1--3.3]{tetra}. Therefore every isometry induces an
automorphism of the checked triangulation. The action is faithful:
fixing each tetrahedron and each of its ideal vertices makes an
isometry the identity. To count automorphisms, choose the image of
one tetrahedron and a permutation of its vertices. Each adjacent
image and vertex permutation is forced by the face pairing.
Propagate through the connected triangulation and reject
inconsistency or failure of bijectivity. The certificate enumerates
all these $24T$ trials on a triangulation of $T$ tetrahedra. The
numbers for $P,N$ are $1,2$ and for $X_1,X_2,X_3$ each is $1$.
This enumeration agrees with the separate SnapPy combinatorial
isomorphism routine. No numerical symmetry guess is a premise.

To check the length bounds, write $m,l$ for the meridian and longitude
translations, with $l$ positive real and area $A=l|\operatorname{Im}m|$.
On cusp zero of $P$, the two quantities
$-\operatorname{Re}m/\sqrt A$ and $3l/\sqrt A$ are strictly larger
than $1$ and $5/2$. On $N$, the quantities
$(-\operatorname{Re}m+2l)/\sqrt A$ and $6l/\sqrt A$ are strictly larger
than $3$ and $5$. The actual normalized length bounds follow by taking
the positive real projection of the slope or its negative. The
verified length-spectrum routine includes the shortest geodesic.
Its lower endpoints exceed $1/2$ for both parents. These are
scale-invariant normalized slope lengths, rather than lengths on the
embedded cusp used for the six-theorem in \cite{covers}.

For $Y_0,Y_1$ the verifier checks the filled hyperbolicity and uses
the verified cutoff enumeration at the stated rational cutoffs. The
length-spectrum contract includes every geodesic below the cutoff
\cite{snappy}. All returned upper endpoints are strictly below that
cutoff. In each case the geodesic labelled with the filling cusp is
identified as the surgery core by the marked geometric structure.
A verified drilling of its word recovers the parent $N$ by an exact
combinatorial isomorphism. This uses the certified geometric drilling
and core detection contract, which identifies the removed geodesic
and its complement.

For $Y_0$, the two geodesic words are $bcckdf$ for the core and $b$
for the other geodesic, in the unsimplified parent presentation.
The exact relators and peripheral words in this presentation are
checked in the same marking as \cite{scope}. The integer homology
images of $b,c,k,d,f$ there are $43,-65,21,24,-19$. Thus the core
has class $43-130+21+24-19=-61$ and the other has class $43$.
The filling imposes $117$ times the marked generator, so the two
classes are distinct even up to sign. The argument needs no assertion
of equality of the two geodesic lengths.
\end{proof}

\begin{lemma}\label{lem:cores}
For $n\ge1$ the isometry group of $Y_n$ has order two and is generated
by the covering involution. For $n=0$, every self-isometry acting
as $-1$ on $H_1(Y_0;\mathbb Z)$ belongs to the subgroup of order two
generated by the covering involution.
\end{lemma}
\begin{proof}
Use the effective filling theorem \cite[Lemma~7.26 and
Theorem~7.28]{fps}. For normalized length at least $L_0=10.1$ it
identifies the single filling core with the unique shortest geodesic
provided the parent systole is larger than
$2\pi/(L_0^2-58)$. Our bounds imply this hypothesis because $\pi<4$
and $8/(L_0^2-58)<1/2$. The theorem applies with the other cusp
unfilled in the exterior case. The normalized length for $X_n$ is
larger than $L_0$ for $n\ge4$, so the unique shortest geodesic drills
to $P$. Lemma~\ref{lem:drilling} implies $\Isom(X_n)=1$ in that range.
For $n=1,2,3$ this conclusion follows directly from the canonical
triangulations in Proposition~\ref{prop:geocert}.

The normalized length for $Y_n$ is larger than $L_0$ for $n\ge2$.
Thus every isometry preserves its core, and Lemma~\ref{lem:drilling}
gives $|\Isom(Y_n)|\le2$. For $Y_1$ the same conclusion comes from
the unique short core in Proposition~\ref{prop:geocert}. The actual
deck transformation has an isometric representative because the
hyperbolic orbifold structure induces the hyperbolic cover metric
\cite{scope}. Rigidity identifies the isometries for this metric and
for the verified filled metric, with their actions on the marked
fundamental group. This supplies the nonidentity isometry in each
case, so the upper bound is attained.

For $Y_0$ any isometry with action $-1$ on homology must preserve
the short core, since the only other geodesic below the cutoff has
an inequivalent class up to sign. The core stabilizer again has order
at most two. The covering involution acts by $-1$ on homology.
Indeed, with four generating orbifold meridians $u_0,\dots,u_3$,
the parity kernel is generated by $u_i u_0$ for $1\le i\le3$.
This follows by grouping any even word in the involutions into pairs.
Conjugation by $u_0$ sends each such generator to its inverse.
Hence the deck transformation belongs to the stabilizer, whose
only elements are the identity and the deck transformation.
\end{proof}

\begin{proof}[Proof of Theorem~\ref{thm:family}]
We have $\rk(\Pi_n)\le3$ by the three products in the preceding
proof. Also $\Pi_n$ is a torsion-free finite-volume lattice, so it
is noncyclic. If it had a generating pair, Lemma~\ref{lem:inversion}
would give an involution inducing a nonidentity isometry of $Y_n$
that acts by $-1$ on $H_1(Y_n;\mathbb Z)$. In every case
Lemma~\ref{lem:cores} identifies this isometry with the deck
transformation. The projective involution would then belong to
$\Gamma_n\setminus\Pi_n$. If it is $J$ and the generating pair is
$(a,b)$, the three elements $J,Ja,Jb$ have order at most two and
generate $\Gamma_n$. This contradicts its torsion-generator rank
four, proving $\rk(\Pi_n)=3$.

Now assume $\Gamma_n$ is generated by two elements. Its abelianization
is $C_2$, as follows from the meridian abelianization of the knot
group. By interchanging the two elements and multiplying one by the
other if necessary, choose a generating pair $(a,b)$ with $a$ odd
and $b$ even for parity. Apply Lemma~\ref{lem:inversion} to the
orbifold lattice. The inversion preserves the characteristic parity
kernel and induces an isometry of $Y_n$. For $n=0$ we need its
homology action. The index-two Schreier generators are
$a^2,b,aba^{-1}$. Inversion sends the first two to their inverses
and the last to $a^{-1}b^{-1}a$, which is conjugate in the kernel
to $(aba^{-1})^{-1}$ by $a^{-2}$. Thus the action on the kernel
abelianization is $-1$, as required. It follows from
Lemma~\ref{lem:cores} that this isometry of $Y_n$ is either the
identity or the deck transformation. In both cases its projective
lift $J$ lies in $\Gamma_n$. Then $J,Ja,Jb$ generate $\Gamma_n$
by at most three torsion elements, a contradiction. Cyclicity is
also impossible for an orbifold lattice of finite covolume. Finally,
$G_n$ surjects onto $\Gamma_n$ and is generated by four meridians.
\end{proof}

\section{The thirty-two named knots}\label{sec:named}
The input knots are as follows, using standard table names.
\begin{center}\small
\begin{tabular}{llll}
$13a_{373}$ & $13a_{375}$ & $13a_{391}$ & $13a_{392}$\\
$13a_{3264}$ & $13a_{3266}$ & $13a_{3275}$ & $13a_{3276}$\\
$13a_{304}$ & $13a_{362}$ & $13a_{623}$ & $13a_{629}$\\
$13a_{3232}$ & $13a_{3259}$ & $13a_{3374}$ & $13a_{3375}$\\
$13n_{834}$ & $13n_{838}$ & $13n_{839}$ & $13n_{991}$\\
$13n_{992}$ & $13n_{993}$ & $13n_{1367}$ & $13n_{1370}$\\
$13n_{1371}$ & $13n_{1374}$ & $13n_{1375}$ & $13n_{1376}$\\
$13n_{3814}$ & $13n_{3840}$ & $13n_{3911}$ & $13n_{3914}$
\end{tabular}
\end{center}
\begin{proof}[Proof of Theorem~\ref{thm:named}]
Use the same marked exterior triangulations as the exact
three-generator proof \cite{uppers}, which binds each to its
actual diagram. The companion computation verifies the complete
hyperbolic structure, strictly negative tilts and canonical
triangulation for each, and the exact tetrahedron-frame enumeration
gives just one automorphism in each case. The isometry groups are
therefore trivial. By Lemma~\ref{lem:inversion}, none of these
noncyclic torsion-free lattices is two-generated. Together with the
three-generator upper bound and meridional equality in \cite{uppers},
this gives the asserted exact ranks.
\end{proof}

The finite geometric computations are specified by
\path{verify_ordinary_rank_from_symmetry.py}. They use verified
arithmetic, canonical geometry and exact enumeration as described
above. The hyperbolic algorithms and the filling theorem are
mathematical inputs. The ordinary ranks of $G_n$ and $\Gamma_n$
within the interval $[3,4]$ remain open.

\paragraph{AI assistance and review.}
ChatGPT (OpenAI) assisted with mathematical arguments, literature checks,
exposition, proof review and verification code.

\begin{thebibliography}{99}
\bibitem{family} \emph{An infinite prime four-bridge family beyond
augmentation and Coxeter bounds}. Companion proof, 18 November 2026.
\bibitem{covers} \emph{Every double branched cover in a prime
four-bridge family is hyperbolic}. Companion proof, 18 November 2026.
\bibitem{scope} \emph{Cyclic branched-cover homology and torsion
generators for a prime four-bridge family}. Companion proof,
18 November 2026.
\bibitem{uppers} \emph{Three-generator upper bounds for thirty-two
knot groups}. Companion proof, 17 November 2026.
\bibitem{bw} M.~Boileau and R.~Weidmann, \emph{The structure of
3-manifolds with two-generated fundamental group}, Topology
44 (2005), 283--320. Author version \href{https://arxiv.org/abs/math/0111327}{arXiv:math/0111327v1}, Lemma~22.
\bibitem{afw} M.~Aschenbrenner, S.~Friedl and H.~Wilton,
\emph{3-Manifold Groups}, EMS, 2015. Rigidity statement in
\href{https://arxiv.org/abs/1205.0202v3}{arXiv:1205.0202v3}.
\bibitem{snappy} M.~Culler, N.~M.~Dunfield, M.~Goerner and
J.~R.~Weeks, SnapPy~3.3.2. \href{https://snappy.computop.org/manifold.html}{Manifold documentation},
verified canonical retriangulation, cusp translations,
length-spectrum and drilling methods.
\bibitem{tetra} E.~Fominykh, S.~Garoufalidis, M.~Goerner,
V.~Tarkaev and A.~Vesnin, \emph{A census of tetrahedral hyperbolic
manifolds}, \href{https://arxiv.org/abs/1502.00383v2}{arXiv:1502.00383v2}, Sections~3.1--3.3.
\bibitem{fps} D.~Futer, J.~S.~Purcell and S.~Schleimer,
\emph{Effective bilipschitz bounds on drilling and filling},
Geometry \& Topology 26 (2022), 1077--1188.
\href{https://doi.org/10.2140/gt.2022.26.1077}{doi:10.2140/gt.2022.26.1077}.
Author version, Lemma~7.26 and Theorem~7.28.
\end{thebibliography}
\end{document}
